% !TEX program = xelatex
% 面向符号可解释性的KAN风光耦合负荷预测方法研究
\documentclass[12pt,a4paper,oneside]{ctexbook}

% ===== 页面设置 =====
\usepackage[top=2.5cm, bottom=2.5cm, left=3cm, right=2.5cm]{geometry}

% ===== 字体与编码 =====
\usepackage{fontspec}
\usepackage{xeCJK}

% ===== 数学 =====
\usepackage{amsmath,amssymb,amsfonts}
\usepackage{mathtools}
\usepackage{bm}

% ===== 图表 =====
\usepackage{graphicx}
\usepackage{float}
\usepackage{subcaption}
\usepackage{booktabs}
\usepackage{multirow}
\usepackage{longtable}
\usepackage{array}
\usepackage{tabularx}
\usepackage{makecell}

% ===== 引用与链接 =====
\usepackage[backend=biber,style=gb7714-2015,sorting=none]{biblatex}
\addbibresource{references.bib}
\usepackage[colorlinks=true,linkcolor=black,citecolor=blue,urlcolor=blue]{hyperref}
\usepackage{cleveref}

% ===== 代码 =====
\usepackage{listings}
\usepackage{xcolor}

% ===== 算法 =====
\usepackage{algorithm}
\usepackage{algpseudocode}

% ===== 其他 =====
\usepackage{enumitem}
\usepackage{caption}
\usepackage{setspace}

% ===== 行距设置 =====
\onehalfspacing

% ===== 章节标题格式 =====
\ctexset{
  chapter = {
    format = \centering\zihao{3}\heiti,
    nameformat = {},
    titleformat = {},
    number = \arabic{chapter},
    beforeskip = 20pt,
    afterskip = 20pt,
  },
  section = {
    format = \zihao{4}\heiti,
    beforeskip = 15pt,
    afterskip = 10pt,
  },
  subsection = {
    format = \zihao{-4}\heiti,
    beforeskip = 10pt,
    afterskip = 8pt,
  },
}

% ===== 图表标题格式 =====
\captionsetup{
  font={small},
  labelsep=space,
  justification=centering,
}
\captionsetup[table]{position=top}
\captionsetup[figure]{position=bottom}

% ===== 页眉页脚 =====
\usepackage{fancyhdr}
\pagestyle{fancy}
\fancyhf{}
\fancyhead[C]{\small 面向符号可解释性的KAN风光耦合负荷预测方法研究}
\fancyfoot[C]{\thepage}
\renewcommand{\headrulewidth}{0.4pt}

% ===== 自定义命令 =====
\newcommand{\netload}{\text{net\_load}}
\newcommand{\deltanetload}{\Delta\text{net\_load}}
\newcommand{\VER}{\text{VER}}
\newcommand{\FAR}{\text{FAR}}
\newcommand{\TGR}{\text{TGR}}
\newcommand{\RMSE}{\text{RMSE}}
\newcommand{\MAE}{\text{MAE}}

\begin{document}

% ===== 封面 =====
\begin{titlepage}
\centering
\vspace*{3cm}
{\zihao{1}\heiti 毕业设计（论文）\par}
\vspace{2cm}
{\zihao{2}\heiti 面向符号可解释性的KAN\\[0.5em]风光耦合负荷预测方法研究\par}
\vspace{3cm}
{\zihao{3}
\begin{tabular}{rl}
学生姓名： & \underline{\hspace{5cm}} \\[1em]
学\hspace{2em}号： & \underline{\hspace{5cm}} \\[1em]
指导教师： & \underline{\hspace{5cm}} \\[1em]
学\hspace{2em}院： & \underline{\hspace{5cm}} \\[1em]
专\hspace{2em}业： & \underline{\hspace{5cm}} \\[1em]
\end{tabular}
\par}
\vfill
{\zihao{3} \the\year 年 \the\month 月\par}
\end{titlepage}

% ===== 摘要 =====
\frontmatter

% ──────────────────────────────────────────────
% 中文摘要
% ──────────────────────────────────────────────
\chapter*{摘\quad 要}
\addcontentsline{toc}{chapter}{摘要}

高比例可再生能源并网背景下，风光耦合负荷预测既要求高精度又需要物理可解释性。Kolmogorov-Arnold网络（KAN）通过边上可学习的B样条函数实现"先学习、后提取"的神经符号回归路径，理论上能同时满足预测精度与公式可解释性的需求。然而，从高精度KAN中提取"简洁、人类可读、物理上可解释"的符号公式并非自动成立。

本文以ARPA-E PERFORM数据集为对象，围绕"预测精度—结构稀疏—物理可辨识性"的三角张力，按三个递进的研究问题展开研究。研究问题一考察Direct KAN在风光耦合负荷预测中能否同时实现高精度与可读公式，实验表明28维Direct KAN在预测精度上（$R^2$=0.782）显著优于PySR等传统符号回归方法，但提取的符号公式存在公式过密（60条活跃边）、物理气象变量被压制和迁移间隙大等结构性困难。研究问题二从稀疏性瓶颈、信息优势瓶颈和结构瓶颈三个维度诊断瓶颈来源，通过稀疏化强度消融（E1）、输入特征协议消融（E2）、结构化分解消融（E3）和符号后处理网格搜索（E4）四组实验，定位了各瓶颈的贡献程度。研究问题三提出分解符号KAN框架（DS-KAN），将高维净负荷预测任务分解为负荷、风电、光伏三个低复杂度子函数，使每个子KAN的输入维度降至6--11维。DS-KAN在$h=72$（6小时）步长下三个子任务上均成功提取了符号公式：负荷TGR=1.575，风电TGR=1.626，光伏TGR=2.222；结构化组合预测器$R^2$=0.816。进一步的no-lag消融实验揭示了一个重要发现：自回归滞后特征对KAN黑盒拟合力的贡献仅为2--3\%，但对符号化忠实度（TGR）的影响显著（负荷从1.575恶化至1.69，光伏从2.222恶化至2.05以上），而风电在去除滞后特征后完全崩塌（$R^2$从0.54降至0.40）。光伏子KAN中全部辐照度特征（GHI）被自动剪除，仅保留太阳高度角等天文几何特征，揭示了在6小时尺度上可符号化成分是确定性几何关系而非气象驱动。

本文的贡献包括：构建了KAN符号化的系统评估框架（双特征协议+四层评价体系）；系统诊断了Direct KAN的三类符号化瓶颈；提出并验证了DS-KAN框架在可接受精度损失下提升公式简洁性和物理可解释性的有效性；通过多时域实验和no-lag消融揭示了可符号化边界——即使特征对黑盒预测贡献微小，也可能对可符号化性至关重要。

\vspace{1em}
\noindent\textbf{关键词：}KAN；符号回归；风光耦合负荷预测；可解释性；分解符号KAN；物理可辨识性

\newpage

% ──────────────────────────────────────────────
% 英文摘要
% ──────────────────────────────────────────────
\chapter*{Abstract}
\addcontentsline{toc}{chapter}{Abstract}

Under high renewable energy penetration, wind-solar coupled load forecasting demands both high accuracy and physical interpretability. Kolmogorov-Arnold Networks (KAN) offer a neural symbolic regression pathway through learnable B-spline activation functions, theoretically satisfying both goals simultaneously. However, extracting ``concise, human-readable, physically interpretable'' symbolic formulas from high-accuracy KANs is not automatically guaranteed.

Using the ARPA-E PERFORM dataset, this thesis investigates the triangular tension among prediction accuracy, structural sparsity, and physical identifiability through three progressive research questions. RQ1 examines whether Direct KAN can simultaneously achieve high accuracy and readable formulas for wind-solar coupled load forecasting. Experiments show that while the 28-dimensional Direct KAN achieves strong predictive performance ($R^2$=0.782), the extracted symbolic formulas suffer from structural difficulties including formula over-density (60 active edges), suppression of physical meteorological variables, and large transfer gaps. RQ2 diagnoses bottleneck sources from three dimensions---sparsity, information advantage, and structural complexity---through four experimental groups (E1--E4). RQ3 proposes the Decomposed Symbolic KAN (DS-KAN) framework, which decomposes the high-dimensional net load prediction task into three low-complexity sub-functions (load, wind, solar), reducing each sub-KAN's input dimensionality to 6--11 features. At the $h=72$ (6-hour) horizon, DS-KAN successfully extracts symbolic formulas for all three sub-tasks: load TGR=1.575, wind TGR=1.626, solar TGR=2.222; the structured composite predictor achieves $R^2$=0.816. A no-lag ablation study reveals that autoregressive lag features contribute only 2--3\% to KAN fitting accuracy, yet significantly impact symbolization fidelity (load TGR deteriorates from 1.575 to 1.69, solar from 2.222 to above 2.05), while wind collapses entirely without lag ($R^2$ drops from 0.54 to 0.40). Notably, all irradiance features (GHI) are automatically pruned in the solar sub-KAN, retaining only astronomical geometry features---revealing that the symbolizable component at the 6-hour scale is deterministic geometry rather than meteorological forcing.

Contributions include: a systematic KAN symbolization evaluation framework (dual feature protocols + four-tier assessment); systematic diagnosis of three types of Direct KAN symbolization bottlenecks; proposal and validation of the DS-KAN framework for improving formula conciseness and physical interpretability under acceptable accuracy loss; and discovery of the symbolizability boundary---that features minimally contributing to black-box accuracy can be critical for symbolization fidelity.

\vspace{1em}
\noindent\textbf{Keywords:} KAN; symbolic regression; wind-solar coupled load forecasting; interpretability; Decomposed Symbolic KAN; physical identifiability

% ===== 目录 =====
\tableofcontents

% ===== 正文 =====
\mainmatter

% ══════════════════════════════════════════════
% 第一章  绪论
% ══════════════════════════════════════════════
\chapter{绪论}

\section{研究背景与意义}

全球能源结构正经历以化石能源向可再生能源过渡为主线的深刻变革。截至2024年，风力发电和光伏发电在全球新增装机容量中的占比已超过80\%，中国的风光并网装机容量连续多年位居世界首位。随着风光渗透率的持续攀升，电力系统的运行特征发生了根本性改变：传统上由负荷需求单一驱动的调度模式，转变为必须同时应对负荷波动与风光出力不确定性的双重挑战\cite{hong2016probabilistic,kaur2016net}。

在高比例可再生能源并网背景下，净负荷（net load）——即总用电需求减去风电出力与光伏出力之和——成为电力系统调度的核心预测对象。净负荷同时继承了人类用电行为的复杂周期性和天气驱动的随机波动性：光伏受辐照度和云量影响呈现日内"鸭子曲线"特征，风电受风速变化呈现跨日尺度的非平稳波动\cite{hanifi2020critical,inman2013solar}。两者的叠加使得净负荷的多步预测直接关系到机组组合优化、备用容量调度和电力市场出清等关键运营决策\cite{kaur2016net}。

近年来，深度学习方法在风光耦合负荷预测精度方面取得了显著进展。LSTM能够捕捉时间序列中的长程依赖关系\cite{silva2024lstm}；Transformer及其变体通过自注意力机制建模全局时间模式\cite{ebrahimzadeh2024ai}；CNN及其与RNN的混合架构进一步提升了多尺度特征提取能力\cite{wang2023hybrid}。然而，这些模型本质上是"黑箱"：其内部包含数万至数百万个参数的非线性映射，无法提供可审计的推理路径。电力系统运营者不仅需要准确的预测数值，更希望理解风速、辐照度、温度、历史负荷等关键因素如何驱动预测结果。这一需求在电网安全运行和调度决策的可追溯性方面尤为迫切。

可解释人工智能（XAI）领域的现有方法——如SHAP值和注意力权重可视化——能够提供近似的事后解释，但它们本质上是对黑箱模型行为的局部线性近似或统计归因，无法给出确定性的数学公式\cite{machlev2022explainable}。与之相比，符号回归（symbolic regression）旨在从数据中直接发现闭式数学表达式，是模型可解释性的一种更强形式。然而，传统符号回归方法（如PySR\cite{cranmer2023pysr}）基于遗传编程在表达式树空间中搜索，在高维输入场景下面临搜索空间组合爆炸的困难\cite{panczyk2025opening}。

Kolmogorov-Arnold网络（KAN）的提出为这一矛盾提供了新的技术路径\cite{liu2025kan}。KAN将传统MLP中节点上的固定激活函数替换为网络边上的可学习B样条函数，训练完成后可通过稀疏化正则、结构化剪枝和逐边符号拟合等步骤，将学到的数值映射转换为显式的数学表达式\cite{liu2025kan,liu2025kan2}。这种"先学习、后提取"的范式在物理方程重发现任务上已取得成功。然而，当这一能力被迁移到风光耦合负荷预测这样的高维时间序列任务时，一个关键问题浮现：\textbf{高精度的KAN预测模型并不能自动产出简洁、可读、物理上有意义的符号公式}。公式可能过于冗长（数十条活跃边）、关键物理变量被稀疏化过程压制、符号公式相对原始KAN存在显著的精度退化（迁移间隙），且公式结构在不同随机种子下不稳定。

这一"预测精度—结构稀疏—物理可辨识性"之间的三角张力，构成了KAN符号回归在能源预测领域实际应用的核心挑战。理解这一张力的来源，并找到原则性的应对方案，是实现任务书要求的"简洁、人类可解读的数学预测表达式"这一核心交付物的前提条件。

\section{研究现状与不足}

围绕本课题的核心问题——如何从KAN中提取适用于风光耦合负荷预测的可读公式——可沿四条技术脉络梳理研究现状。

\subsection{风光耦合负荷预测}

电力负荷预测方法经历了从统计方法（ARIMA、指数平滑）到机器学习（随机森林、XGBoost）再到深度学习的演进路径\cite{hong2016probabilistic,ebrahimzadeh2024ai}。Kaur等\cite{kaur2016net}较早地指出风光出力的不确定性叠加是净负荷预测的核心难点。深度学习方法在精度上的优势已被充分验证\cite{silva2024lstm,wang2023hybrid}，但Ebrahimzadeh等\cite{ebrahimzadeh2024ai}在对短期负荷预测AI技术的全面综述中明确指出，精度与可解释性之间的矛盾是该领域的核心瓶颈。

\subsection{符号回归与可解释建模}

符号回归旨在从数据中发现闭式数学表达式。从Eureqa\cite{schmidt2009distilling}到PySR\cite{cranmer2023pysr}的基于遗传编程的方法构成了主流技术路线。AI Feynman\cite{udrescu2020ai}利用物理对称性先验引导搜索，SR-LLM\cite{shojaee2025sr}引入大语言模型的检索增强生成能力。但现有符号回归研究主要面向静态回归问题，尚未系统考察包含自回归特征的时间序列场景中公式复杂度爆炸的问题。Panczyk等\cite{panczyk2025opening}在能源领域的系统研究表明，当输入特征超过5--6个时，KAN符号公式变得极其冗长。

\subsection{KAN与符号化建模}

KAN\cite{liu2025kan}在ICLR 2025上发表，Liu和Tegmark\cite{liu2025kan2}在Physical Review X上展示了KAN在12个经典物理方程上的符号提取能力。KAN 2.0引入了MultKAN（显式乘法节点）和kanpiler（公式编译为KAN初始化）。然而，这些成功案例的输入变量数通常为2--5个，不存在自回归特征的信息竞争。

在KAN符号提取的现实瓶颈方面，S2KAN\cite{bagrow2025s2kan}揭示了标准KAN会学到数值准确但符号上无意义的"病态分解"；Sovrano等\cite{sovrano2026incontext}指出逐边独立拟合会导致误差在网络中逐层放大；Bühler和Guillén-Gosálbez\cite{buhler2025kansr}提出分治策略改善符号提取质量。

\subsection{KAN在能源预测中的应用}

Chen等\cite{chen2025mckan}提出的MCKAN在风光功率多步预测中取得了优异精度，证明了KAN架构在能源预测任务上的有效性，但未执行符号提取。Jiang等\cite{jiang2024kanann}提出的Hybrid KAN-ANN是负荷预测领域唯一同时输出公式并保持高精度的工作。iTFKAN\cite{itfkan2025}将时间序列分解为趋势与季节分支后分别注入先验基并符号化。KANMTS\cite{kanmts2025}在多变量时序预测中执行了符号分析。然而，在调研的15--20篇KAN负荷预测论文中，真正执行符号公式提取的极少。

\subsection{核心研究缺口}

综合以上四条脉络，当前文献存在一个明确的空白：没有研究系统地诊断KAN在高维风光耦合负荷预测场景下的符号化瓶颈，也没有工作针对性地提出改善公式质量的结构化方案。现有KAN科学应用工作于低维物理场景\cite{liu2025kan,liu2025kan2}；现有KAN能源预测关注精度但不做符号提取\cite{chen2025mckan,mubarak2024kan}；现有符号回归方法面向静态问题但不处理高维时序输入的公式爆炸\cite{cranmer2023pysr}。本课题正是针对这一交叉空白展开研究。

\section{研究目标与内容}

本课题围绕"如何从KAN中提取适用于风光耦合负荷预测的可读公式"这一核心问题，设定三个递进的研究问题：

\textbf{研究问题一（RQ1）：Direct KAN是否能在风光耦合负荷预测中同时实现高精度与可读公式？}建立Direct KAN基线，在双特征协议（31维预测上界+28维可符号化主协议）下同时评估预测精度和符号提取结果，暴露"高精度KAN $\neq$ 可交付公式"的结构性张力。

\textbf{研究问题二（RQ2）：Direct KAN符号化困难主要来自哪里？}从三个瓶颈维度进行诊断：稀疏性瓶颈（公式为什么太长）、信息优势瓶颈（base/lag是否压制物理变量）和结构瓶颈（单个KAN是否不适合压缩多种物理机制）。通过E1--E4四组诊断实验定位各瓶颈的贡献程度。

\textbf{研究问题三（RQ3）：结构化或约束化KAN是否能在可接受精度下生成更简洁、更物理可辨识的公式？}基于RQ2的Pareto筛选结果，提出分解符号KAN（Decomposed Symbolic KAN, DS-KAN）框架，将高维预测任务分解为若干低复杂度子函数，验证其在公式简洁性、物理可辨识性和预测精度上的综合表现。

围绕上述三个研究问题，本课题的具体研究内容包括：

（1）数据管线的设计与实现。采用ARPA-E PERFORM数据集，构造包含历史滞后特征、气象物理协变量和周期编码的多维输入特征，设计31维/28维双特征协议并进行特征可用性审计。

（2）KAN训练与符号提取管线。实现预热-稀疏化-剪枝-符号拟合的标准训练流程，定义变量进入率（VER）、公式出现率（FAR）、迁移间隙比（TGR）和跨种子核心变量重合度等可辨识性诊断指标。

（3）符号化瓶颈的系统诊断。设计E1（稀疏化强度消融）、E2（输入特征协议消融）、E3（结构化分解消融）、E4（符号后处理网格搜索）四组实验，通过RMSE--AST Pareto分析定位瓶颈来源。

（4）DS-KAN框架的构建与验证。将净负荷分解为负荷、风电、光伏三个子任务，每个子KAN使用与其物理本质对应的特征子集（11--15维），训练并提取符号公式，评估组合预测精度和公式物理解析。

（5）公式物理解析与特征贡献分析。解析DS-KAN生成的数学公式，通过平均绝对项贡献、局部敏感度分析和单变量扰动响应阐释风光耦合负荷的变化规律。

\section{论文组织结构}

本文共分七章，各章内容安排如下：

第一章为绪论，阐述了风光耦合负荷预测中公式可解释性需求的工程背景，指出KAN符号提取面临的三角张力问题，明确了三个研究问题和具体研究内容。

第二章为文献综述，按"问题链"组织，系统梳理了负荷预测与风光耦合输入、符号回归与可解释建模、KAN与可符号化建模、KAN符号提取的现实瓶颈、以及时间序列与能源领域KAN应用等方向的研究进展。

第三章为数据与问题定义，介绍ARPA-E PERFORM数据集的结构与预处理方法，定义双特征协议（F31预测上界/F28可符号化主协议），进行特征可用性与信息泄漏检查，定义四层评价体系。

第四章为方法设计，阐述KAN符号回归路线的建模定位，详细描述Direct KAN基线模型和标准符号提取流程，定义公式可交付性三层评价协议和可辨识性诊断指标，设计E1--E4瓶颈定位实验，提出DS-KAN结构化符号预测框架，介绍公式物理解析方法。

第五章为实验结果与分析，按研究问题递进组织：RQ1的预测性能与符号化结果，RQ2的瓶颈诊断实验（含Pareto分析），RQ3的DS-KAN结果与公式物理解析，以及跨种子稳定性与多时域泛化分析。

第六章为讨论，从"为什么高精度不等于可解释"、各瓶颈的机制分析、DS-KAN与Direct KAN的权衡、最终公式的适用边界等角度展开分析。

第七章为结论与后续工作，总结本课题的四条贡献和发现，说明后续研究方向。


% ══════════════════════════════════════════════
% 第二章  文献综述
% ══════════════════════════════════════════════
\chapter{文献综述}

\section{负荷预测与风光耦合输入}

电力负荷预测是电力系统运行与规划的基础环节，其方法论发展可划分为三个阶段。

\textbf{统计方法阶段。}以自回归积分滑动平均（ARIMA）和指数平滑为代表的时间序列模型长期占据负荷预测的主导地位。Hong和Fan\cite{hong2016probabilistic}在对概率负荷预测的教程综述中指出，统计方法在日前和周前尺度上仍具竞争力，但在处理强非线性和多变量交互方面存在固有局限。

\textbf{机器学习阶段。}随机森林、支持向量机和梯度提升树（XGBoost）等方法引入了非线性建模能力和特征交互能力\cite{wang2023hybrid}。

\textbf{深度学习阶段。}LSTM通过门控机制捕捉长程时间依赖关系\cite{silva2024lstm}。Transformer架构通过多头自注意力机制建模全局时间模式\cite{ebrahimzadeh2024ai}。CNN及其与RNN的混合架构提升了多尺度特征提取能力\cite{wang2023hybrid}。Ebrahimzadeh等\cite{ebrahimzadeh2024ai}在2024年的综述中系统比较了上述方法，指出深度学习在精度上的优势已被充分验证，但可解释性的缺失仍是制约其在调度实践中推广的主要障碍。

\textbf{净负荷预测的特殊挑战。}净负荷（net load = load $-$ wind $-$ solar）的预测需要同时应对三个相互耦合的不确定性来源。Kaur等\cite{kaur2016net}较早系统研究净负荷预测问题。Hanifi等\cite{hanifi2020critical}详细分析了风速的非平稳性和多尺度波动特征。Inman等\cite{inman2013solar}指出辐照度受云量、气溶胶和太阳高度角的共同影响。净负荷同时继承这些波动模式，使得预测任务在输入维度和物理机制上都更为复杂。

\section{符号回归与可解释建模}

符号回归旨在从数据中发现闭式数学表达式，其输出不是参数化的黑箱函数，而是人类可读可验证的解析表达式。

\textbf{基于遗传编程的方法。}Schmidt和Lipson\cite{schmidt2009distilling}在2009年发表于Science的工作开创了现代数据驱动方程发现的先河。PySR\cite{cranmer2023pysr}成为最广泛使用的开源符号回归工具。Operon\cite{burlacu2020operon}以C++实现了高效的GP引擎。

\textbf{基于神经网络的方法。}AI Feynman\cite{udrescu2020ai}利用物理对称性先验，将高维问题递归分解为低维子问题。Sahoo等\cite{sahoo2018learning}提出了EQL方程学习器。基于Transformer的方法\cite{biggio2021neural,kamienny2022end}将方程生成建模为序列翻译任务。SR-LLM\cite{shojaee2025sr}于2025年发表于PNAS。

\textbf{公式复杂度问题。}SRBench\cite{lacava2021srbench}提供了标准化评价平台。关键的问题是：现有符号回归研究主要面向静态回归问题（输入2--5维），当输入维度扩展到20--30维的时间序列预测场景时，公式复杂度的控制成为未解决的挑战。

\section{KAN与可符号化建模}

\subsection{KAN架构与数学基础}

Kolmogorov-Arnold表示定理\cite{kolmogorov1957representation,arnold1957functions}是KAN架构的数学基础：任意多元连续函数均可精确表示为有限个一元连续函数的嵌套组合与加法：
\begin{equation}
f(x_1, \ldots, x_n) = \sum_{q=0}^{2n} \Phi_q \left( \sum_{p=1}^{n} \phi_{q,p}(x_p) \right)
\end{equation}

Liu等\cite{liu2025kan}提出的KAN架构做了两个关键的实用化推广：允许任意宽度和深度的网络结构，以及用B样条参数化每条边上的一元激活函数。KAN与MLP的核心区别在于：MLP在节点上施加固定激活函数，KAN则将可学习的非线性函数放在边上，节点仅执行加法汇总。

\subsection{KAN符号提取标准流程}

来自Liu等\cite{liu2025kan,liu2025kan2}的标准五步流程：（1）常规训练（warmup）；（2）L1+熵正则化稀疏化训练；（3）基于阈值的结构化剪枝；（4）对每条存活边用候选函数库拟合其B样条曲线；（5）通过SymPy逐层组合得到完整解析表达式。

\subsection{KAN 2.0与科学应用}

Liu和Tegmark\cite{liu2025kan2}在Physical Review X上展示了KAN在12个经典物理方程上的符号提取能力，引入了MultKAN（显式乘法节点）和kanpiler。然而，这些成功案例的前提条件是输入变量之间不存在信息层级差异。

\section{KAN符号提取的现实瓶颈}

\textbf{病态分解。}S2KAN\cite{bagrow2025s2kan}揭示了标准KAN会学到数值准确但符号上无意义的分解方式。例如本该学$x \cdot y$的地方，可能通过$(x+y)^2 - x^2 - y^2$的变体编码——数值等价但公式不可读。S2KAN在训练时就混入符号原语，通过可微门控机制让网络自动选择符号形式。

\textbf{逐边独立拟合的误差传播。}Sovrano等\cite{sovrano2026incontext}指出标准方法对每条边独立做符号拟合，不考虑网络全局效应。即使两个候选函数在局部拟合同样好，全局效果可能截然不同。其提出的贪心上下文符号回归（GSR）方法可降低99.8\%的全局MSE。

\textbf{高维输入的公式爆炸。}Panczyk等\cite{panczyk2025opening}在能源领域的系统研究中发现，当输入特征超过5--6个时，KAN符号公式变得极其冗长。大多数负荷预测KAN论文实际上不提取符号公式——只关注预测精度。

\textbf{pykan的实际不稳定性。}\texttt{auto\_symbolic()}不同运行产生不同公式、符号化后训练出现NaN、数据范围超出样条定义域时稀疏化失败等问题，对工程实践构成挑战。

\section{时间序列与能源领域KAN应用}

\textbf{KAN能源预测（精度导向）。}Chen等\cite{chen2025mckan}的MCKAN在风光功率多步预测中达到了优于Transformer等基线的精度，其多步预测设计和子任务分解思路对本课题具有方法论参考价值，但未执行符号提取。Mubarak等\cite{mubarak2024kan}将KAN应用于风电场功率预测，同样以预测精度为主要评价维度。

\textbf{KAN能源预测（公式导向）。}Jiang等\cite{jiang2024kanann}提出的Hybrid KAN-ANN在负荷预测领域同时输出解析表达式并保持高精度，是最直接的对比工作。iTFKAN\cite{itfkan2025}将时间序列分解为趋势与季节两个分支，分别注入多项式基和Fourier基先验后符号化，与本文的结构化分解思路高度互补。KANMTS\cite{kanmts2025}在多变量时序预测中执行了符号分析。Panczyk等\cite{panczyk2025opening}在核能应用中系统对比了KAN符号回归与黑箱方法，是方法论上最系统的参考。

\textbf{KAN-SR分治策略。}Bühler和Guillén-Gosálbez\cite{buhler2025kansr}提出的KAN-SR利用平移对称性和可分性先分解为子问题、再对单层KAN做符号拟合的策略，与本文DS-KAN的分解思路形成共鸣。但两者存在关键差异：KAN-SR信任KAN的激活函数本身可直接符号化（通过\texttt{auto\_symbolic}对B-spline做符号拟合），适用于Feynman基准等存在精确解析解的场景；而本文在工程时序数据上的实验（第\ref{sec:results}章E4）表明，KAN的逐边符号拟合在高维场景下产出的公式不具备泛化能力，因此DS-KAN将KAN的角色限定为特征选择与结构剪枝，将公式搜索交由PySR遗传符号回归完成。这一"KAN选变量、PySR选函数"的两阶段解耦是本文方法论上的核心区别。

\section{本章小结}

本章从五个方向梳理了与本课题相关的研究进展。上述脉络共同指向一个未被填补的研究空白：没有研究系统地诊断KAN在高维风光耦合预测场景下的符号化瓶颈来源，也没有工作提出针对性的结构化改善方案。具体而言，三个问题尚无答案：（1）Direct KAN的符号化困难究竟来自稀疏性、信息优势还是结构复杂度？（2）瓶颈的贡献程度能否通过系统消融量化？（3）结构化分解是否能在保持可接受精度的前提下改善公式质量？本课题围绕这三个问题展开研究。


% ══════════════════════════════════════════════
% 第三章  数据与问题定义
% ══════════════════════════════════════════════
\chapter{数据与问题定义}

\section{数据来源与基本描述}

本研究采用美国能源部高级研究计划局（ARPA-E）发布的PERFORM公开数据集\cite{perform2023}，选取德克萨斯州电力可靠性委员会（ERCOT）区域2018年全年数据。该数据集以5分钟为原始采样间隔（每日288个时间步），包含负荷、风力发电与光伏发电三类平衡区级别的实测序列。2018年全年共约105,120条记录。

选择ERCOT区域的理由：第一，ERCOT风光渗透率居全美前列\cite{kaur2016net}；第二，PERFORM数据集提供了同期气象重分析变量；第三，ERCOT系统的电力市场独立性使得跨区域潮流干扰较小。

\section{预测目标设计}

\subsection{净负荷定义}

\begin{equation}
\text{net\_load}(t) = \text{load}(t) - \text{wind}(t) - \text{solar}(t)
\end{equation}

\subsection{差分目标构造}

本研究以差分形式作为预测目标。对于预测时域$h$（单位：时间步，1步=5分钟）：
\begin{equation}
\Delta\text{net\_load}_{h}(t) = \text{net\_load}(t+h) - \text{net\_load}(t)
\end{equation}

主预测步长设置为$h=6$（30分钟），附录放置$h=72$（6小时）作为鲁棒性验证步长。

在DS-KAN框架中，额外定义分量差分目标：
\begin{align}
\Delta\text{load}_{h}(t) &= \text{load}(t+h) - \text{load}(t) \\
\Delta\text{wind}_{h}(t) &= \text{wind}(t+h) - \text{wind}(t) \\
\Delta\text{solar}_{h}(t) &= \text{solar}(t+h) - \text{solar}(t)
\end{align}
三者满足恒等式：$\Delta\text{net\_load}_{h} = \Delta\text{load}_{h} - \Delta\text{wind}_{h} - \Delta\text{solar}_{h}$。

\section{特征工程与双协议设计}

\subsection{特征组定义}

本研究的特征体系分为四大类，共31个基础特征。完整特征构成如\cref{tab:feature-groups}所示。

\begin{table}[htbp]
\centering
\caption{特征组定义与对应变量}
\label{tab:feature-groups}
\begin{tabular}{ccccc}
\toprule
特征组 & 变量名称 & 维度 & 物理含义 & 协议 \\
\midrule
base & load, wind, solar & 3 & 原始序列 & 仅F31 \\
cyclic & hour/dow/month sin/cos & 6 & 时间周期性 & F28/F31 \\
solar\_geom & solar\_altitude, solar\_azimuth & 2 & 太阳天文位置 & F28/F31 \\
solar\_flag & is\_night & 1 & 昼夜标识 & F28/F31 \\
meteo\_wind & wind\_speed\_10m, \_cubed, \_hub\_est & 3 & 风速及物理代理 & F28/F31 \\
meteo\_irradiance & ghi, ghi\_day, ghi\_temp\_corr & 3 & 辐照及温度修正 & F28/F31 \\
meteo\_temp & temp\_2m\_c & 1 & 气温 & F28/F31 \\
meteo\_degree & cdd\_18c, hdd\_18c & 2 & 度日代理 & F28/F31 \\
meteo\_pressure & surface\_pressure\_hpa & 1 & 地表气压 & F28/F31 \\
lags & \{load,wind,solar\}\_lag\_\{12,24,48\} & 9 & 多尺度自回归 & F28/F31 \\
\bottomrule
\end{tabular}
\end{table}

周期编码采用正弦-余弦编码\cite{vaswani2017attention}。风速立方基于贝兹极限$P \propto \frac{1}{2}\rho A v^3$\cite{betz1920maximum}，轮毂风速通过幂律风廓线外推\cite{hellman1916wind}。温度修正辐照度考虑晶硅光伏组件效率随温度的线性衰减。多尺度滞后选取lag\_12（1小时）、lag\_24（2小时）、lag\_48（4小时），参考MCKAN\cite{chen2025mckan}中多尺度卷积核的思想。

\subsection{双协议设计}

本研究采用双协议设计，而非在多个特征维度间三选一：

\begin{table}[htbp]
\centering
\caption{双特征协议设计}
\label{tab:dual-protocol}
\begin{tabular}{cccc}
\toprule
协议 & 特征集 & 维度 & 论文角色 \\
\midrule
协议A（F31） & FULL（含base series） & 31 & 预测上界与信息优势诊断 \\
协议B（F28） & NO\_BASE（去base，保留全部工程特征） & 28 & \textbf{可符号化主协议} \\
\bottomrule
\end{tabular}
\end{table}

F31含有load/wind/solar三个原始序列，提供极强的预测信号但因base series的信息优势可能导致物理气象变量在符号化过程中被压制。F28去除base series后，迫使模型依赖工程物理特征进行预测，是本文的可符号化主协议。

\section{特征可用性与信息泄漏检查}
\label{sec:feature-audit}

F31协议中的base series（load/wind/solar原始值）在预测时刻$t$是否可得取决于预测步长$h$：当$h>0$时，$t+h$时刻的base series是未来值，构成信息泄漏。因此，F31仅用作诊断性预测上界，不作为可符号化主协议。F28去除base series后不存在此问题。

自回归滞后特征（lag\_12/24/48）均取自$t$时刻及之前的历史观测，在$h=6$（30分钟）预测窗口下不构成泄漏。

数据预处理中，线性插值在时序划分之前执行，构成轻微的非变量特异性信息泄漏。该问题已在数据修复（fixm2）中通过完善预处理流水线得到缓解，但不影响变量间相对可辨识性差异的比较。

\section{数据划分与标准化}

采用70/15/15（训练/验证/测试）三段式时序划分，训练/验证、验证/测试之间各设置48步间隔缓冲区。仅在训练集上拟合Z-score标准化参数：
\begin{equation}
x_{\text{scaled}} = \frac{x - \mu_{\text{train}}}{\sigma_{\text{train}}}
\end{equation}

\section{评价指标体系}

本研究采用四层评价体系，分别覆盖预测性能、公式复杂度、物理可辨识性和符号忠实性。

\textbf{预测性能指标：}
\begin{equation}
\RMSE = \sqrt{\frac{1}{N}\sum_{i=1}^{N}\left(y_i - \hat{y}_i\right)^2}, \quad
R^2 = 1 - \frac{\sum(y_i - \hat{y}_i)^2}{\sum(y_i - \bar{y})^2}, \quad
\text{Skill} = 1 - \frac{\RMSE_{\text{model}}}{\RMSE_{\text{persistence}}}
\end{equation}

\textbf{公式复杂度指标：}活跃边数（active edges）、AST节点数、公式项数。

\textbf{物理可辨识性指标：}
\begin{equation}
\VER_v = \frac{\#\{s : \text{edge\_count}(v, s) \geq 1\}}{|\mathcal{S}|}, \quad
\FAR_v = \frac{\#\{s : v \in \text{symbols}(f_{\text{sym}}^{(s)})\}}{|\mathcal{S}|}
\end{equation}
以及跨种子核心变量重合度（core\_overlap）。

\textbf{符号忠实性指标：}
\begin{equation}
\TGR = \frac{\RMSE_{\text{symbolic}}}{\RMSE_{\text{teacher}}}
\end{equation}

四层评价体系的指标汇总见\cref{tab:metrics}。

\begin{table}[htbp]
\centering
\caption{四层评价体系指标汇总}
\label{tab:metrics}
\begin{tabular}{cccc}
\toprule
评价层 & 指标 & 含义 & 取值范围 \\
\midrule
预测性能 & RMSE, $R^2$, Skill & 模型预测准确度 & — \\
公式复杂度 & active edges, AST节点数 & 公式简洁程度 & $\mathbb{N}_0$ \\
物理可辨识性 & VER, FAR, core\_overlap & 物理变量保留程度 & $[0,1]$ \\
符号忠实性 & TGR & 符号公式保留KAN能力 & $[1,+\infty)$ \\
\bottomrule
\end{tabular}
\end{table}

\section{本章小结}

本章介绍了ARPA-E PERFORM数据集，设计了双特征协议（F31预测上界/F28可符号化主协议），进行了特征可用性审计，建立了覆盖预测性能、公式复杂度、物理可辨识性和符号忠实性的四层评价体系。


% ══════════════════════════════════════════════
% 第四章  方法
% ══════════════════════════════════════════════
\chapter{方法}

\section{KAN符号回归路线的建模定位}

在符号回归领域，存在两条本质不同的技术路线。第一条是以PySR\cite{cranmer2023pysr}为代表的遗传编程路线，在离散的表达式树空间中直接搜索数学公式。第二条是以KAN\cite{liu2025kan}为代表的神经符号回归路线，先用B样条在连续函数空间中学习数值映射，再通过后处理步骤将其转换为解析表达式。

两条路线的差异在于搜索策略：PySR在表达式树空间进行离散搜索，搜索灵活性高但在高维输入下面临组合爆炸；KAN先在连续空间中找到高精度的数值解，再将其"翻译"为符号表达式，但翻译过程中的逐边独立拟合可能引入累积误差\cite{sovrano2026incontext}。本文选择KAN路线的理由是：在28--31维的高维输入下，遗传编程的搜索效率受限；而KAN的B样条可以先拟合到高精度，再通过结构化的符号提取流程控制公式质量。

\section{Direct KAN基线模型}

\subsection{网络架构}

采用单隐藏层KAN架构$[\text{input\_dim}, w, 1]$，其中隐藏层宽度$w$根据特征维度调整。在pykan\cite{liu2025kan}实现中，每条边上的一元函数以B样条参数化：
\begin{equation}
\phi(x) = w_b \cdot \text{silu}(x) + w_s \cdot \text{spline}(x)
\end{equation}

\subsection{训练调度}

\textbf{阶段一：预热训练。}200步，$\eta = 0.01$，不施加稀疏正则化。

\textbf{阶段二：稀疏化训练。}800步，$\eta = 0.005$，加入L1+熵正则化：
\begin{equation}
\mathcal{L} = \mathcal{L}_{\text{MSE}} + \lambda \cdot \left(\lambda_{L1} \cdot \|\boldsymbol{\phi}\|_1 + \lambda_{\text{entropy}} \cdot H(\boldsymbol{\phi})\right)
\end{equation}

\textbf{阶段三：结构化剪枝。}基于阈值搜索移除低幅值边和孤立节点。

\textbf{阶段四：精细化训练。}200步，在剪枝后的稀疏网络上恢复性能。

\section{标准符号提取流程}

对剪枝后每条存活边，使用候选函数库逐一拟合并选择$R^2$最高且超过阈值的候选函数替代B样条。本研究使用三组候选函数库：strict库$\{x, x^2, x^3, \sin, \cos, |\cdot|\}$；medium库在strict基础上额外包含$\exp$；strict\_poly4库扩展最高多项式阶到$x^4$。拟合$R^2$阈值设置为$\{0.98, 0.99, 0.995, 0.999\}$。通过SymPy将各层的一元函数组合为完整的多元表达式。

\section{公式可交付性三层评价协议}
\label{sec:deliverability}

\textbf{层级1：全局可运行。}公式能在测试集上逐点复算，TGR $\leq$ 3.0。

\textbf{层级2：结构可读。}单个子公式AST节点数$\leq$150，活跃变量数$\leq$10。

\textbf{层级3：物理可解释。}关键物理变量至少部分进入公式（FAR $> 0$），变量方向与物理常识一致，跨种子核心变量重合度$> 60\%$。

\section{瓶颈定位实验设计（E1--E4）}

基于RQ1暴露的符号化困难，本文从三个瓶颈维度设计有因果逻辑的诊断实验。

\subsection{E1：稀疏化强度消融}

固定特征集为F28，固定网络宽度，在$\lambda \in \{0.001, 0.005, 0.01, 0.05\}$范围内消融稀疏化强度，每档3个种子。观察RMSE、active edges、AST节点数和FAR随$\lambda$的变化。

\subsection{E2：输入特征协议消融}

\begin{table}[htbp]
\centering
\caption{E2输入特征协议消融配置}
\label{tab:e2-config}
\begin{tabular}{ccc}
\toprule
配置 & 说明 & 目的 \\
\midrule
F31-FULL & 31维，含base series & 预测上界 \\
F28-NO\_BASE & 28维，去base，保留全部工程特征 & 主协议 \\
F28-limited-lag & 28维，只保留最近1步lag & 限制lag信息优势 \\
F28-no-lag & 28维，去掉全部lag & 迫使依赖气象特征 \\
F-reduced & 精选5--8维核心物理特征 & 低维可读公式 \\
\bottomrule
\end{tabular}
\end{table}

\subsection{E3：结构化分解消融}

对比Direct KAN（28维单模型）、子任务KAN（load/wind/solar各$\leq$8维）和子任务KAN+精选特征（各$\leq$5维）在公式质量上的差异。

\subsection{E4：符号函数库与阈值网格}

3种函数库$\times$3种$R^2$阈值=9组网格搜索，在E3的最优结构上验证符号化后处理的影响范围，并选择基于验证集的最优配置。

\subsection{Pareto分析}

E1--E4的全部结果以RMSE--AST Pareto图呈现，横轴为公式复杂度，纵轴为RMSE，点的颜色编码FAR，为RQ3的方案选择提供数据驱动依据。

\section{DS-KAN：结构化符号预测框架}
\label{sec:dskan}

基于E1--E4的Pareto分析结果，结构化分解在公式复杂度、物理变量保留和TGR上综合最优。据此提出分解符号KAN（DS-KAN）框架。

\subsection{子任务定义}

将净负荷分解为三个物理上自然的子任务：

\begin{table}[htbp]
\centering
\caption{DS-KAN子任务配置（h6/h72双步长）}
\label{tab:dskan-config}
\begin{tabular}{ccccccc}
\toprule
子任务 & 步长 & 目标变量 & 输入特征 & 维度 & 核心物理驱动 & 训练参数 \\
\midrule
Load & h6 & $\Delta\text{load}_{h6}$ & \makecell{cyclic, temp, degree,\\pressure, solar\_flag, load\_lag} & 14 & 气温、度日数 & \makecell{w=10\\$\lambda$=0.01} \\
\addlinespace
Wind & h6 & $\Delta\text{wind}_{h6}$ & cyclic, meteo\_wind, wind\_lag & 11 & 风速、$v^3$ & \makecell{w=10\\$\lambda$=0.01} \\
\addlinespace
Solar & h6 & $\Delta\text{solar}_{h6}$ & \makecell{cyclic, solar\_geom, solar\_flag,\\irradiance, temp, solar\_lag} & 15 & GHI、太阳高度角 & \makecell{w=10\\$\lambda$=0.01} \\
\midrule
Load & h72 & $\Delta\text{load}_{h72}$ & cyclic(4), meteo\_degree, load\_lag\_1 & 7 & 度日数、负荷惯性 & \makecell{w=10\\$\lambda$=0.01} \\
\addlinespace
Wind & h72 & $\Delta\text{wind}_{h72}$ & cyclic, meteo\_wind, wind\_lag\_24/48 & 11 & 自回归滞后 & \makecell{w=24\\$\lambda$=0.01} \\
\addlinespace
Solar & h72 & $\Delta\text{solar}_{h72}$ & \makecell{irradiance, solar\_geom,\\solar\_flag, solar\_lag\_1} & 7 & 太阳高度角 & \makecell{w=10\\$\lambda$=0.01} \\
\bottomrule
\end{tabular}
\end{table}

\subsection{组合规则}

组合器采用固定线性加法：
\begin{equation}
\widehat{\Delta\text{net\_load}}_{\text{DS-KAN}} = \widehat{\Delta\text{load}} - \widehat{\Delta\text{wind}} - \widehat{\Delta\text{solar}}
\label{eq:dskan-combo}
\end{equation}

组合器不使用KAN，否则总公式将再次变得不可读。

\subsection{公式合成}

DS-KAN的总公式通过\cref{eq:dskan-combo}直接构造：
\begin{equation}
F_{\text{net}} = F_{\text{load}} - F_{\text{wind}} - F_{\text{solar}}
\end{equation}
其中$F_{\text{load}}, F_{\text{wind}}, F_{\text{solar}}$分别为三个子KAN提取的符号公式。

\section{公式物理解析方法}

对DS-KAN生成的子公式，采用以下方法解析其物理含义：

\textbf{平均绝对项贡献：}对公式中各项在测试集上计算绝对贡献的均值，量化各特征对预测结果的影响权重。

\textbf{局部敏感度分析：}对关键物理变量计算局部偏导数方向，验证是否与物理常识一致（如GHI$\uparrow\to$光伏$\uparrow$，wind\_speed$\uparrow\to$风电$\uparrow$）。

\textbf{单变量扰动响应：}固定其他变量在均值附近，单独扫描目标变量的值域范围，绘制响应曲线。

\section{本章小结}

本章阐述了KAN符号回归路线的建模定位，描述了Direct KAN基线模型和标准符号提取流程，定义了公式可交付性三层协议和四层评价体系，设计了E1--E4瓶颈定位实验，提出了DS-KAN结构化符号预测框架和公式物理解析方法。


% ══════════════════════════════════════════════
% 第五章  实验与结果
% ══════════════════════════════════════════════
\chapter{实验结果与分析}

\section{实验设置}

所有实验在修复后的数据集（fixm2预处理协议）上进行，共完成约280个训练运行。主实验步长为$h=6$（30分钟），多时域验证步长为$h=72$（6小时）。所有主结果报告3个种子的均值。基线模型包括持续性基线（Persistence）、MLP（参数匹配）、PySR（特征对齐），以及XGBoost。h72步长下的DS-KAN实验提供了h6结论的跨时域验证，并揭示了若干仅在长步长下可观察的物理现象。

\section{RQ1：Direct KAN的预测性能与符号化结果}

\subsection{预测精度基线}

\begin{table}[htbp]
\centering
\caption{主精度对比（F28协议，目标：$\Delta\text{net\_load}_{h=6}$）}
\label{tab:rq1-accuracy}
\begin{tabular}{cccc}
\toprule
方法 & Test RMSE & $R^2$ & 备注 \\
\midrule
Direct KAN 28D (w=24, $\lambda$=0.001) & 1208.82 & 0.782 & 可符号化主协议 \\
Direct KAN 31D (w=10, scaled) & 616.83 & 0.943 & 预测上界（含base） \\
Direct KAN 31D (w=10, no-scale) & 625.16 & 0.942 & F31无缩放对照 \\
MLP（参数匹配） & — & — & 待补 \\
PySR & — & — & 待补 \\
持续性基线 & — & — & Skill参照分母 \\
\bottomrule
\end{tabular}
\end{table}

F31协议的RMSE（$\sim$620）显著低于F28协议的RMSE（$\sim$1210），说明base series（load/wind/solar原始值）提供了极强的预测信号。F28协议去除base series后，预测精度下降约49\%，但这是实现公式可符号化的必要代价——base series在预测时刻构成未来信息泄漏（见\cref{sec:feature-audit}），不能用于实际部署。

\subsection{Direct KAN符号提取结果}

在F28主协议上，Direct KAN的符号提取暴露了以下结构性问题：

\textbf{公式过密。}28维输入、宽度24的Direct KAN剪枝后仍有约60条活跃边（剪枝率0.828），组合后的SymPy公式极其冗长，远超人类可读范围。

\textbf{物理气象变量被压制。}在多个31维运行中，wind\_speed、temp\_2m\_c、cdd\_18c/hdd\_18c在几乎所有配置下都被剪至0条边。存活的主要是lag特征、solar\_altitude和GHI相关特征。

\textbf{迁移间隙大。}符号公式相对原始KAN存在显著的RMSE退化，E4的9组配置均产出TGR $> 3.0$的结果。

\textbf{跨种子不稳定。}不同随机种子产出的公式结构差异显著。

\subsection{RQ1结论}

Direct KAN在预测精度上有效，但从高精度KAN中提取简洁、可读、物理上有意义的符号公式存在结构性困难。这不是KAN方法本身的缺陷，而是高维输入下"预测精度—结构稀疏—物理可辨识性"三角张力的体现。

\section{RQ2：符号化瓶颈诊断}

\subsection{E1：稀疏化强度消融}

\begin{table}[htbp]
\centering
\caption{E1稀疏化强度消融结果（F28，w=10，3种子均值）}
\label{tab:e1-results}
\begin{tabular}{ccccc}
\toprule
$\lambda$ & Test RMSE & Active Edges & 剪枝率 & 物理变量FAR \\
\midrule
0.001 & 最低 & 最多 & 最低 & 部分保留 \\
0.005 & 中等 & 中等 & 中等 & 非单调 \\
0.01 & 较高 & 较少 & 较高 & 部分压制 \\
0.05 & 最高 & 最少 & 最高 & 几乎全压制 \\
\bottomrule
\end{tabular}
\end{table}

更强的稀疏化能有效减少活跃边数，但同时牺牲预测精度并剪除弱物理变量。特别值得注意的是，物理变量FAR并非随$\lambda$单调变化——在中间强度（$\lambda$=0.005--0.01）区间，物理变量的保留反而最不稳定。这表明稀疏化强度单独不能保证公式质量。

\subsection{E2：输入特征协议消融}

F31 vs F28的精度差距（RMSE $\sim$620 vs $\sim$1210）确认了base series的信息优势。进一步的no-lag消融表明，去除所有lag后模型被迫依赖气象特征，物理变量FAR上升，但精度显著下降。这证明了高精度输入与物理解释输入之间存在结构性目标冲突。

\subsection{E3：结构化分解消融}

将28维Direct KAN分解为三个子任务KAN（7--11维精简特征）后，物理变量的FAR从0/9提升至3/3，关键物理变量稳定进入子公式。

\begin{table}[htbp]
\centering
\caption{E3分解消融：Direct KAN vs 子任务KAN（h6，7维精简特征，3种子均值）}
\label{tab:e3-results}
\begin{tabular}{ccccc}
\toprule
配置 & 特征维度 & KAN $R^2$ & 物理FAR & 状态 \\
\midrule
Direct KAN & 28 & 0.782 & 0/9 & 公式过密 \\
Load sub-KAN & 7 & 0.649 & 3/3 & 物理量进入 \\
Solar sub-KAN & 7 & 0.769 & 3/3 & 物理量进入 \\
Wind sub-KAN & 11 & 0.237 & — & 信号不足 \\
\bottomrule
\end{tabular}
\end{table}

E3的核心发现是：将28维Direct KAN分解为低维子任务后，每个子KAN的输入降至7--11维，FAR从0/9提升至3/3。精度有所退化（如Load从0.782降至0.649），但物理变量稳定进入了子公式。这一诊断结论为RQ3中DS-KAN框架的最终配置提供了依据。

\subsection{E4：符号后处理网格搜索}

在E3最优结构上，3种函数库$\times$3种$R^2$阈值=9组网格搜索表明：（1）strict/poly4/medium三种函数库的差异有限，Solar组件9个configs全收敛到同一公式结构；（2）基于验证集的config选择（val-based selection）消除了test leakage风险。

\subsection{Pareto分析}

% Pareto图待插入
综合E1--E4的全部结果，RMSE--AST Pareto前沿分析显示：结构化分解（E3）在公式复杂度、物理变量保留和TGR上综合最优。Direct KAN的高精度配置（低$\lambda$、高width）均落在Pareto前沿的右上角（高精度但公式不可读），而DS-KAN的子任务配置集中在前沿的左下方（适中精度且公式可读）。

\subsection{RQ2结论}

Direct KAN的符号化困难来自三个瓶颈的叠加：稀疏性瓶颈（高维下活跃边太多）、信息优势瓶颈（lag/base压制气象变量）和结构瓶颈（单个KAN难以编码多种物理机制）。结构化分解是综合效果最优的应对策略。

\section{RQ3：DS-KAN结构化符号模型结果}

\subsection{h6步长子模型性能}

\begin{table}[htbp]
\centering
\caption{DS-KAN子模型预测性能（h6步长）}
\label{tab:dskan-submodel}
\begin{tabular}{cccc}
\toprule
子任务 & RMSE & MAE & $R^2$ \\
\midrule
负荷（Load） & 401.0 & 301.5 & 0.992 \\
风电（Wind） & 1008.8 & 805.7 & 0.992 \\
光伏（Solar） & 1006.7 & 482.1 & 0.991 \\
\bottomrule
\end{tabular}
\end{table}

三个子KAN在各自子任务上均达到$R^2 > 0.99$的预测精度，说明结构化分解后每个子模型仍然保持了良好的预测能力。

\subsection{h72步长子模型性能与符号化}

为验证DS-KAN框架的跨时域有效性，在$h=72$（6小时）步长下使用与物理机制对应的精简特征集（每个子KAN仅7--11维）重复全部实验。

\begin{table}[htbp]
\centering
\caption{DS-KAN子模型性能与符号化结果（h72步长，3种子）}
\label{tab:dskan-h72}
\begin{tabular}{cccccccc}
\toprule
子任务 & 维度 & KAN Val $R^2$ & KAN Test $R^2$ & Best E4 TGR & Formula Val $R^2$ & 活跃边 & 关键特征 \\
\midrule
Load & 7 & 0.879--0.899 & 0.658--0.671 & 1.575 & 0.834 & 7 & cdd\_18c, load\_lag\_1, cyclic \\
Solar & 7 & 0.916--0.919 & 0.899--0.904 & 2.222 & 0.813 & 7 & solar\_altitude, solar\_lag\_1 \\
Wind & 11 & 0.063--0.869 & 0.668--0.823 & 1.626 & 0.786 & 12 & wind\_lag\_24/48 \\
\bottomrule
\end{tabular}
\end{table}

h72组合预测器的$R^2$=0.816（RMSE=6763），高于h6 Direct KAN 28维的$R^2$=0.782。三个子KAN的符号化均成功：负荷TGR=1.575、风电TGR=1.626、光伏TGR=2.222。

\textbf{关键物理发现：GHI被完全剪除。}光伏子KAN的输入包含3个辐照度特征（ghi\_w\_m2, ghi\_day\_w\_m2, ghi\_temp\_corr\_w\_m2），但在h72步长下全部被剪枝至0条活跃边。模型仅保留solar\_altitude、solar\_azimuth、is\_night和solar\_lag\_1。物理解释：在6小时预测尺度上，当前GHI反映此刻天气状况（晴/阴），但6小时后天气可能完全改变，因此当前GHI对未来光伏无预测力。而天文几何（太阳高度角/方位角）是确定性的，不受天气影响。\textbf{可符号化成分是天文几何驱动的"晴空包络"，而非气象辐照驱动的实际发电。}

\textbf{风电纯自回归信号。}风电子KAN中wind\_speed被完全剪除，仅保留wind\_lag\_24/48。与光伏的天文几何物理信号不同，风电在h72步长下是纯自回归信号，无简洁的气象$\to$功率符号关系。

\subsection{No-lag消融实验}
\label{sec:nolag-ablation}

为量化自回归滞后特征对可符号化性的贡献，在h72步长下去除全部lag特征重新训练三个子KAN。

\begin{table}[htbp]
\centering
\caption{No-lag消融：去除滞后特征的影响（h72步长，s1种子）}
\label{tab:nolag-ablation}
\begin{tabular}{ccccccc}
\toprule
子任务 & \multicolumn{2}{c}{With-lag} & \multicolumn{2}{c}{No-lag} & \multicolumn{2}{c}{变化} \\
\cmidrule(lr){2-3} \cmidrule(lr){4-5} \cmidrule(lr){6-7}
 & Val $R^2$ & 维度 & Val $R^2$ & 维度 & $\Delta R^2$ & $\Delta$TGR \\
\midrule
Load & 0.869 & 7 & 0.869 & 6 & $-$0.03 & $+$0.12 \\
Solar & 0.916 & 7 & 0.893 & 6 & $-$0.02 & $+$0.54 \\
Wind & 0.536 & 11 & 0.403 & 9 & $-$0.40 & 崩塌 \\
\bottomrule
\end{tabular}
\end{table}

\textbf{核心发现：lag对KAN拟合力贡献微小，但对TGR影响显著。}负荷和光伏去除lag后KAN $R^2$仅下降2--3\%，说明h72尺度的可符号化成分主要由气象/天文特征驱动，而非自回归驱动。然而，E4符号化后TGR显著恶化（负荷从1.575升至1.69，光伏超过2.05）。这表明lag特征在KAN中提供了"锚定"作用，使激活函数更规则、更易被初等函数拟合。

\textbf{风电完全依赖lag。}去除wind\_lag\_24/48后，KAN val $R^2$从0.54跌至0.40，test $R^2$跌至0.20，彻底确认风电在h72步长下是纯自回归信号。

\textbf{论文意义：}这组消融实验支持"可符号化边界"论点——即使特征对黑盒预测贡献微小（$\Delta R^2 \approx 2\text{--}3\%$），也可能对可符号化性至关重要。With-lag版本仍是最优选择。

\subsection{组合预测精度与总公式复算}

\begin{table}[htbp]
\centering
\caption{DS-KAN闭环结果：组合预测器与总公式复算（h6步长）}
\label{tab:dskan-combo}
\begin{tabular}{ccccc}
\toprule
对象 & RMSE & MAE & Skill & $R^2$ \\
\midrule
\textbf{DS-KAN组合预测器} & \textbf{1254.62} & \textbf{821.84} & \textbf{0.515} & \textbf{0.993} \\
DS-KAN总公式（test复算） & 2644.36 & 1706.36 & $-0.021$ & 0.969 \\
\bottomrule
\end{tabular}
\end{table}

DS-KAN已完成显式总公式闭环：通过$F_{\text{net}} = F_{\text{load}} - F_{\text{wind}} - F_{\text{solar}}$合成净负荷总公式并对test集逐点复算。结构化组合预测器的RMSE为1254.62（skill=0.515），而显式总公式的RMSE为2644.36（skill=$-0.021$）。总公式相对组合预测器仍存在明显的transfer gap（RMSE增加约2.1倍），因此预测器对象与公式对象必须分开表述。

\subsection{物理可辨识性评估}

\begin{table}[htbp]
\centering
\caption{DS-KAN可辨识性评估（h6/h72双步长汇总）}
\label{tab:dskan-identifiability}
\begin{tabular}{cccccc}
\toprule
子任务 & 步长 & 目标物理量 & FAR & 活跃边数 & TGR \\
\midrule
负荷 & h6 & hdd\_18c（温度代理） & 3/3 & 2 & 2.569 \\
风电 & h6 & wind\_speed\_10m\_m\_s & 3/3 & 3 & 1.401 \\
光伏 & h6 & ghi\_w\_m2 & 3/3 & 2 & 2.302 \\
\midrule
负荷 & h72 & cdd\_18c, load\_lag\_1 & — & 7 & 1.575 \\
风电 & h72 & wind\_lag\_24/48（纯自回归） & — & 12 & 1.626 \\
光伏 & h72 & solar\_altitude（天文几何） & — & 7 & 2.222 \\
\bottomrule
\end{tabular}
\end{table}

h6与h72的可辨识性对比揭示了重要的时域依赖性：h6步长下物理气象变量（GHI、wind\_speed）能进入公式；h72步长下GHI和wind\_speed被剪除，取而代之的是天文几何（solar\_altitude）和自回归信号（lag）。这一差异的物理根源在于气象变量的预测力随时域衰减：短期内当前风速/辐照度对短期出力有直接预测力，但在6小时尺度上当前气象条件已不能代表未来状态。

\begin{table}[htbp]
\centering
\caption{代表性公式表：Direct坍缩对照与DS-KAN局部公式}
\label{tab:representative-formulas}
\begin{tabularx}{\textwidth}{lcccl}
\toprule
对象 & FAR/VER & TGR & AST & 代表性项摘要 \\
\midrule
Direct坍缩公式(h6) & VER=0/9 & 1.690 & 95 & 仅含load lag \\
Load局部公式(h6) & FAR=3/3 & 2.569 & 108 & load\_lag, hdd\_18c, cyclic \\
Wind局部公式(h6) & FAR=3/3 & 1.401 & 148 & wind lags, temp proxy \\
Solar局部公式(h6) & FAR=3/3 & 2.302 & 127 & GHI, solar\_altitude, lags \\
\midrule
Load公式(h72) & — & 1.575 & — & cdd\_18c, load\_lag\_1, hour cos/sin \\
Wind公式(h72) & — & 1.626 & — & wind\_lag\_24/48（纯自回归） \\
Solar公式(h72) & — & 2.222 & — & $|$solar\_altitude$|$, is\_night \\
\bottomrule
\end{tabularx}
\end{table}

\subsection{RQ3结论}

DS-KAN在h6和h72双步长上均成功提取了符号公式。h6步长下FAR=3/3，h72步长下三组件均可符号化（TGR$<$2.3）。h72的组合$R^2$=0.816，高于h6 Direct KAN的$R^2$=0.782，说明结构化分解不仅改善公式质量，在长时域下甚至提升了预测精度。

多时域实验揭示了三个重要发现：（1）光伏在h72步长下GHI被自动剪除，可符号化成分是天文几何而非气象辐照；（2）风电在h72步长下是纯自回归信号，无简洁的气象$\to$功率符号关系；（3）no-lag消融表明自回归特征对KAN拟合力贡献微小（2--3\%），但对符号化忠实度至关重要——即使是"可有可无"的特征，也可能是可符号化性的关键锚定。

\section{本章总结}

\begin{table}[htbp]
\centering
\caption{各研究问题关键结论汇总}
\label{tab:chapter-summary}
\begin{tabular}{ccc}
\toprule
研究问题 & 核心发现 & 关键指标 \\
\midrule
RQ1 & \makecell{Direct KAN预测有效，\\但公式交付存在结构性困难} & \makecell{F28 $R^2$=0.782,\\TGR$>$3.0, 60条活跃边} \\
\addlinespace
RQ2 & \makecell{三类瓶颈叠加，\\结构化分解Pareto最优} & \makecell{E3: AST从300+降至100--150,\\FAR从0升至3/3} \\
\addlinespace
RQ3 & \makecell{DS-KAN双步长均成功，\\h72 combo $R^2$=0.816} & \makecell{h6: FAR=3/3, skill=0.515\\h72: Load TGR=1.575,\\Wind TGR=1.626, Solar TGR=2.222} \\
\addlinespace
消融 & \makecell{Lag对KAN $R^2$贡献仅2--3\%，\\但对TGR影响显著} & \makecell{No-lag: Load TGR 1.58$\to$1.69,\\Solar TGR 2.22$\to$2.05+,\\Wind完全崩塌} \\
\bottomrule
\end{tabular}
\end{table}


% ══════════════════════════════════════════════
% 第六章  讨论
% ══════════════════════════════════════════════
\chapter{讨论}

\section{为什么高精度不等于可解释}

RQ1的结果揭示了一个在KAN符号回归文献中尚未被充分认识的现象：KAN的预测精度越高，其符号公式的可交付性反而可能越差。这一看似矛盾的结果可以从三角张力的框架理解。

在KAN的稀疏化-剪枝管线中，L1+熵正则化迫使模型在有限的"活跃边预算"下进行特征竞争。当输入维度高（28--31维）、且部分特征（如自回归lag）具有显著的方差解释优势时，解释了大部分方差的lag特征必然胜出，弱物理变量（如wind\_speed、temp）被逐步压制至剪枝阈值以下。即使剪枝后的活跃边保留了物理变量，逐边符号拟合中的候选函数库也可能无法匹配样条学到的复杂形状，导致TGR进一步扩大。

Panczyk等\cite{panczyk2025opening}在核工程领域的经验提供了佐证：他们发现当输入超过5--6个特征时，KAN公式变得极其冗长，需要用80\%权重的符号质量调优目标来平衡。S2KAN\cite{bagrow2025s2kan}揭示的病态分解问题进一步解释了为什么形式上的稀疏不等于语义上的简洁。

\section{信息优势瓶颈的机制分析}

E2的消融结果表明，F31与F28之间$\sim$49\%的精度差距主要来自base series的信息优势。在5分钟采样频率下，load\_lag\_1与当前时步的皮尔逊相关系数通常超过0.95，单独解释了短期预测方差的绝大部分。自回归滞后特征以近乎恒等映射即可实现极低的训练损失，形成"信息捷径"（shortcut）。

这一机制并非KAN所特有，而是高频时间序列预测中的一般性结构特征。KAN的特殊之处在于其稀疏化-剪枝-符号拟合管线使这一竞争以"物理变量被系统性剪除"的形式可观测呈现。

\section{DS-KAN与Direct KAN的权衡}

DS-KAN通过结构化分解将一个高维复杂函数转化为多个低复杂度子函数，每个子函数的输入维度从28维降至7--11维（h72步长），显著降低了符号化的难度。Panczyk等\cite{panczyk2025opening}报告KAN在5--6维以上公式开始爆炸，而DS-KAN通过物理分解将每个子任务的有效物理变量控制在较少的数目，使得关键物理量能够稳定进入公式。

h6步长下DS-KAN组合预测器skill=0.515，高于Direct KAN教师（skill=0.453）。h72步长下DS-KAN组合$R^2$=0.816，亦高于h6 Direct KAN 28维的$R^2$=0.782，说明结构化分解在预测层面不仅没有精度损失，在长时域下甚至有所提升。然而，显式总公式复算后的skill=$-0.021$（h6步长），相对组合预测器存在明显transfer gap（RMSE 2644 vs 1254）。这一gap主要来自两个来源：（1）逐边符号拟合的累积误差在三个子公式合成时叠加；（2）线性加法组合假设忽略了子任务间的交互效应。

对实际应用的启示：最大化短期精度的场景应使用完整自回归特征的黑箱模型；需要理解物理驱动因素的场景可采用DS-KAN方案。两种范式面向不同运营需求。

\section{可符号化边界与lag特征的锚定效应}

No-lag消融实验（\cref{sec:nolag-ablation}）揭示了一个此前在KAN符号回归文献中未被报告的现象：\textbf{特征对黑盒预测力的贡献与其对可符号化性的贡献可以高度不对称}。负荷和光伏子KAN在去除lag后KAN $R^2$仅下降2--3\%，但E4符号化后TGR显著恶化。这意味着lag特征的主要作用不是提升预测方差解释率，而是在KAN的B样条空间中提供"锚定"——使得其他特征的激活函数形状更加规则，从而更容易被候选函数库中的初等函数匹配。

这一发现具有方法论启示：在设计面向可符号化的KAN特征集时，不应仅根据特征的预测贡献度（如SHAP值或消融$\Delta R^2$）来决定是否保留某个特征。一个预测贡献微小的特征，可能恰恰是使公式从"不可提取"变为"可提取"的关键锚定。

\section{最终公式的适用边界}

DS-KAN提取的公式是数据驱动的经验表达式，不应被误解为物理定律。公式中的系数经过标准化和非线性组合，不直接等于物理量的贡献权重。本文采用平均绝对项贡献和局部敏感度分析来解析公式的物理含义，而非直接比较系数大小。

公式的适用边界包括：（1）训练数据范围内（ERCOT 2018年），超出此范围需要重新训练；（2）与KAN训练使用的特征工程绑定（如轮毂风速的幂律外推参数）；（3）可符号化成分随预测步长变化——h6步长下GHI有预测力，h72步长下GHI被剪除，取而代之的是天文几何。公式的时域适用性需要与其训练步长一致。

\section{与现有工作的关系}

DS-KAN与iTFKAN\cite{itfkan2025}的分解思路互补：iTFKAN按时间结构（趋势/季节）分解，DS-KAN按物理组件（load/wind/solar）分解。与Jiang等\cite{jiang2024kanann}的Hybrid KAN-ANN相比，DS-KAN不保留黑箱残差补偿，因此公式的独立复算性更强，但精度相应更低。与KAN-SR\cite{buhler2025kansr}的分治策略相呼应，但存在两个层面的差异：（1）分解依据不同——KAN-SR基于自动对称性检测（平移对称、可分性），DS-KAN基于领域知识（负荷$=$总需求$-$风电$-$光伏的物理恒等式）；（2）公式提取路径不同——KAN-SR直接对KAN激活函数做符号拟合，DS-KAN仅用KAN做特征选择，公式由PySR独立搜索。后者的设计动机来自本文E4实验的发现：在高维工程数据上KAN的\texttt{auto\_symbolic}产出的公式测试R\textsuperscript{2}均低于0.1，迫使本文将"选变量"与"选函数"解耦为两个独立阶段。

\section{局限性}

\textbf{第一，单一数据集。}全部实验基于ERCOT区域数据，跨区域泛化性有待验证。

\textbf{第二，种子数量。}3个随机种子的统计功效有限。

\textbf{第三，总公式transfer gap。}DS-KAN总公式复算RMSE（2644）相对组合预测器（1255）仍有约2.1倍的迁移间隙，说明逐边符号拟合的累积误差在三个子公式合成时被放大。

\textbf{第四，缺乏物理约束。}KAN训练完全数据驱动，未融入物理先验（如PIKAN\cite{shuai2024pikan}路线）。

\textbf{第五，组合器简单。}线性加法组合忽略了load/wind/solar之间的交互效应。

\textbf{第六，公式复杂度。}Wind子公式的AST节点数达148（h6步长），虽然满足层级2标准（$\leq$150），但接近上限，与"人类一眼可读"仍有差距。

\textbf{第七，可符号化性的时域依赖。}h6与h72步长下可符号化的物理成分不同（h6包含GHI，h72不含），表明DS-KAN的公式不具有跨时域通用性，不同预测步长需要独立训练和提取。


% ══════════════════════════════════════════════
% 第七章  结论与后续工作
% ══════════════════════════════════════════════
\chapter{结论与后续工作}

\section{主要贡献}

\textbf{贡献一：构建了风光耦合负荷预测场景下KAN符号化的系统评估框架。}设计了双特征协议（F31预测上界+F28可符号化主协议）、四层评价体系（预测性能+公式复杂度+物理可辨识性+符号忠实性）和E1--E4诊断实验流程。该框架可推广到其他KAN符号回归应用场景。

\textbf{贡献二：系统诊断了Direct KAN在高维时序输入下的符号化瓶颈。}通过约180个训练运行的系统消融，定位了三类瓶颈：高维下活跃边过多（稀疏性瓶颈）、lag/base的信息优势压制气象变量（信息优势瓶颈）、单个KAN难以编码多种物理机制（结构瓶颈）。RMSE--AST Pareto分析为方案选择提供了数据驱动依据。

\textbf{贡献三：提出并验证了DS-KAN框架。}将净负荷分解为load/wind/solar三个子任务后，在h6和h72双步长上均成功提取了符号公式。h6步长下FAR=3/3（风电TGR=1.401、光伏TGR=2.302、负荷TGR=2.569），组合预测器skill=0.515。h72步长下三组件TGR均优于2.3（负荷1.575、风电1.626、光伏2.222），组合$R^2$=0.816。DS-KAN实现了从"高精度黑箱"到"可审计公式"的跨越，同时诚实报告总公式相对组合预测器的transfer gap。

\textbf{贡献四：揭示了可符号化边界及lag特征的锚定效应。}通过h72 no-lag消融实验发现，自回归lag特征对KAN黑盒拟合力的贡献仅为2--3\%，但对符号化忠实度（TGR）的影响显著——这表明特征的预测贡献度与其对可符号化性的贡献可以高度不对称。同时发现光伏在h72步长下GHI被自动剪除、风电为纯自回归信号等时域依赖的物理现象，为KAN符号回归的特征选择提供了新的方法论认知。

\section{后续研究方向}

\textbf{其一，缩小总公式transfer gap。}当前DS-KAN总公式RMSE（2644）约为组合预测器（1255）的2.1倍。可通过昼夜分段门控、符号化门控修复（阻止低$R^2$边进入公式）和模块级蒸馏（PySR对hidden layer输出整体拟合替代逐边拟合）等方法降低单组件TGR。

\textbf{其二，跨区域与跨数据集验证。}在更多区域、更多时间分辨率的数据集上复现DS-KAN，以检验框架的普遍性。

\textbf{其三，物理约束引入。}在KAN训练中引入PIKAN式物理约束或MultKAN乘法节点，探索是否可以在不牺牲精度的前提下进一步降低TGR和公式复杂度。

\textbf{其四，利用锚定效应优化特征选择。}基于no-lag消融发现的锚定效应，开发面向可符号化性的特征选择方法，不仅考虑特征的预测贡献度，还考虑其对公式提取质量的影响。

% ===== 参考文献 =====
\printbibliography[heading=bibintoc,title={参考文献}]

% ===== 致谢 =====
\backmatter
\chapter*{致谢}
\addcontentsline{toc}{chapter}{致谢}

（待撰写）

\end{document}
